\ifdefined\HMTAPPFormal
\chapter{APP : support, opérations et chemins}
\label{chap:app}
\index{APP}

\noindent
Le sigle APP provient de \emph{All Possible Paths}. Il désigne une unique
structure de chemins sur un support arithmétique fini : le graphe
sous-jacent possède quatre-vingt-un sommets, tandis que l'ensemble des chemins
finis, réuni sur toutes les longueurs, est dénombrable. Sa matrice
arithmétique principale est la table multiplicative nonadique. L'addition est la loi
interne dont l'accumulation engendre chaque ligne et chaque colonne de cette table et agit
aussi comme observable de comparaison sur le même support. APP réunit ainsi
addition, produit et
composition de trajectoires avant d'effectuer la projection ternaire ou
l'évaluation archimédienne. La définition suivante fixe de manière exhaustive
l'univers de chemins considéré.

\section{Domaine, générateurs et relations préformelles du support APP}

Le support d'APP est un tore arithmétique fini : chacune de ses deux
coordonnées parcourt les neuf classes résiduelles et revient à son origine après
neuf translations unitaires. Ses quatre-vingt-un points sont donc des
paires de classes modulo neuf. Le graphe de Cayley associé conserve cette
structure cyclique : un sommet représente une position, une arête orientée
représente un déplacement cardinal unitaire et une suite finie
d'arêtes représente un chemin. Ainsi, \emph{tous les chemins
possibles} désigne l'ensemble des mots cardinaux finis ayant leur point
initial dans le tore.

Chaque sommet reçoit simultanément deux évaluations arithmétiques. Le
\emph{feuillet additif} publie la somme de ses deux coordonnées et le
\emph{feuillet multiplicatif} publie leur produit. Le terme \emph{feuillet}
désigne ici une évaluation du même point et du même chemin ; APP conserve
un unique support et deux fonctions sur celui-ci.

La différence opératoire entre les deux feuillets est exacte. Dans une ligne ou une colonne,
l'addition agit par translation et accumule de manière répétée le même incrément.
La multiplication par une unité permute également les classes ; la
multiplication par un diviseur de zéro replie plusieurs classes sur un même
résidu. Pour conserver l'information que ce repliement masque, APP élève
chaque évaluation à ses coordonnées résidu--quotient :
\[
 i+j=R^+_{ij}+9q^+_{ij},
 \qquad
 ij=R^\times_{ij}+9q^\times_{ij}.
\]
Les paires \((R^+,q^+)\) et \((R^\times,q^\times)\) récupèrent les valeurs
entières antérieures à la réduction. Lors de la concaténation de blocs, la retenue additive
\(\kappa_9\) enregistre l'écart entre additionner des représentants entiers et
additionner d'abord leurs résidus ; avec les deux quotients, elle rend la
récupération compositionnelle. Ainsi, l'accumulation additive construit les trajectoires visibles,
tandis que le relèvement résidu--quotient déploie les fibres créées par le
repliement multiplicatif.

Pour un chemin \(\gamma=(x_0,\ldots,x_r)\), la généalogie visible est la
suite des positions et des paires d'évaluations
\[
 \mathcal G_{\rm vis}(\gamma)
 =\bigl(\gamma;(\Sigma(x_k),\Pi(x_k))_{k=0}^{r}\bigr).
\]
Sa généalogie enrichie ajoute, à chaque pas, les quotients, la retenue et
l'orientation transportée. La projection enrichie--visible conserve le
chemin et les résidus et oublie précisément ces coordonnées de mémoire. Cette
distinction sera l'entrée de l'état complet du TRIT et du Topological Phase Kernel.
Le domaine complet des (104\,976) conditions initiales à double curseur et
les (468) fibres d'émission qu'elles produisent sont publiés intégralement dans
\cref{app:catalogo-semillas-app-rev9} ; chaque germe y est distingué de
l'émission à laquelle il appartient et la multiplicité de la fibre est conservée.

\section{Support arithmétique, translations et observables}

Soit
\[
\dr(n)=1+((n-1)\bmod9)\in\{1,\ldots,9\}.
\]
Cette application choisit le représentant \(9\) pour la classe nulle. Nous définissons
\(X_9=(\ZZ/9\ZZ)^2\), sur lequel agit le groupe de translations
\((\ZZ/9\ZZ)^2\), et fixons l'ensemble symétrique de générateurs
\[
 \mathscr D=\{(1,0),(0,1),(-1,0),(0,-1)\}.
\]
Le graphe de Cayley
\[
 \GraphAPP=\operatorname{Cay}(X_9,\mathscr D)
\]
possède \(81\) sommets et quatre arêtes orientées ayant leur origine à chaque sommet.
Un chemin orienté de longueur \(r\) est une suite
\(\gamma=(x_0,\ldots,x_r)\) dans \(X_9\) telle que
\(x_{k+1}-x_k\in\mathscr D\) pour tout \(k<r\). Nous notons
\(\operatorname{Path}(\GraphAPP)\) l'ensemble de tous les chemins finis
ainsi définis.

\begin{definition}[Structure APP]
La structure arithmétique des chemins possibles est le tuple
\[
 \AAPP=(X_9,\GraphAPP,\Sigma,\Pi,
        \operatorname{Path}(\GraphAPP)),
\]
où les deux observables de sommet sont
\[
\Sigma(i,j)=\dr(i+j),
\qquad
\Pi(i,j)=\dr(ij),
\qquad1\le i,j\le9.
\]
\end{definition}

\subsection*{Cadre cardinal et mots de trajectoire}

Pour rendre explicite l'orientation du tore, nous utilisons les coordonnées
\((\text{ligne},\text{colonne})\), avec les lignes croissantes vers le sud et les
colonnes croissantes vers l'est, et écrivons
\[
 \delta(N)=(-1,0),\qquad \delta(E)=(0,1),\qquad
 \delta(S)=(1,0),\qquad \delta(O)=(0,-1).
\]
Soit \(\mathscr D_{\rm card}=\{E,N,O,S\}\), où \(O\) désigne l'ouest. Étant donnés un point initial
\(x_0\in X_9\) et un mot \(w=d_1\cdots d_r\in
\mathscr D_{\rm card}^{r}\), sa trajectoire est
\[
 \gamma(x_0;w)_k
 =x_0+\sum_{j=1}^{k}\delta(d_j)\pmod 9,
 \qquad 0\leq k\leq r.
\]
La correspondance \((x_0,w)\mapsto\gamma(x_0;w)\) est bijective entre
\(X_9\times\mathscr D_{\rm card}^{r}\) et les chemins orientés de
longueur \(r\). En particulier, il existe exactement \(81\cdot4^r\) chemins
de cette longueur lorsque les retours et les reculs sont autorisés.

Le mot \(w\) se ferme sur le tore précisément lorsque
\[
\#S(w)-\#N(w)\equiv0\pmod9,
\qquad
\#E(w)-\#O(w)\equiv0\pmod9.
\]
Pour un mot fermé, le déplacement entier calculé avant réduction
modulo neuf définit le vecteur d'enroulement
\[
 \operatorname{wind}(w)=\frac1{9}
 \bigl(\#S(w)-\#N(w),\#E(w)-\#O(w)\bigr)\in\ZZ^2.
\]
L'inversion de l'orientation inverse l'ordre du mot, échange
\(E\leftrightarrow O\) et \(N\leftrightarrow S\), et transforme
\(\operatorname{wind}(w)\) en son opposé. Ce cadre cardinal est celui qui est
transporté ultérieurement dans le noyau de phase et fait partie de la définition
de chaque trajectoire.

Les expressions sont écrites ici dans la carte des représentants
\(\{1,\ldots,9\}\). La formule
\(\widetilde\Sigma(i,j)=\dr(i+j-1)\) est une carte translatée. Elle représente le
même tore uniquement lorsque la translation est également appliquée aux indices, aux conditions
initiales, aux trajectoires positionnelles et aux étiquettes. Toutes les démonstrations suivantes emploient
\(\Sigma(i,j)=\dr(i+j)\) et n'alternent pas entre les deux cartes.

\begin{figure}[H]
\centering
\resizebox{\textwidth}{!}{\begin{tikzpicture}[font=\scriptsize,cell/.style={minimum width=4.2mm,minimum height=4.2mm,inner sep=0pt,draw=hmtgray!35}]
  \begin{scope}[xshift=0cm]
    \node[font=\bfseries\scriptsize,hmtblue,align=center] at (1.72,.92)
      {matrice arithmétique APP\\[-.5mm]\(\Pi(i,j)=\dr(ij)\)};
    \foreach \i in {1,...,9}{
      \node[font=\bfseries] at (-.35,{-(\i-1)*.43}) {\i};
      \foreach \j in {1,...,9}{
        \pgfmathtruncatemacro{\v}{mod(\i*\j-1,9)+1}
        \pgfmathtruncatemacro{\border}{(\i==9)||(\j==9)}
        \pgfmathtruncatemacro{\rad}{(mod(\i,3)==0)||(mod(\j,3)==0)}
        \ifnum\border=1
          \node[cell,fill=hmtlightblue,text=hmtblue,font=\bfseries] at ({(\j-1)*.43},{-(\i-1)*.43}) {\v};
        \else
          \ifnum\rad=1
            \node[cell,fill=hmtlightviolet] at ({(\j-1)*.43},{-(\i-1)*.43}) {\v};
          \else
            \node[cell,fill=white] at ({(\j-1)*.43},{-(\i-1)*.43}) {\v};
          \fi
        \fi
      }
    }
    \foreach \j in {1,...,9}{\node[font=\bfseries] at ({(\j-1)*.43},.38) {\j};}
    \draw[hmtgold,very thick,rounded corners=1pt] (1.075,-1.075) rectangle (1.935,-1.935);
    \draw[hmtblue,very thick] (3.72,.22)--(3.72,-3.72)--(-.22,-3.72);
    \node[hmtblue,anchor=north] at (1.75,-4.02) {subconjunto coordenado de la clase nula};
  \end{scope}

  \begin{scope}[xshift=5.25cm]
    \node[draw=hmtblue,rounded corners=2pt,fill=hmtlightblue,
      text width=6.25cm,align=left,inner sep=8pt] at (2.9,-1.35) {%
      \textbf{Ley interna de acumulación}\par\smallskip
      \[
        \Pi(i,j)=\dr(\underbrace{i+\cdots+i}_{j\ {\rm sumandos}})
        =\dr(ij).
      \]
      La somme engendre les lignes multiplicatives et, dans la direction
      transversale, les colonnes. L'observable
      \(\Sigma(i,j)=\dr(i+j)\) conserve cette loi élémentaire pour comparer
      les deux lectures sur le même point.};
    \node[draw=hmtgold,rounded corners=2pt,fill=hmtlightgold,
      text width=6.25cm,align=center,inner sep=8pt] at (2.9,-3.65) {%
      \(\displaystyle
      \AAPP=(X_9,\GraphAPP,\Pi;\Sigma,
      \operatorname{Path}(\GraphAPP))\)\par
      une seule APP, une matrice principale, une loi additive interne et quatre
      directions orientées \(N,E,S,O\).};
  \end{scope}
\end{tikzpicture}
}
\caption{L'unique matrice APP et sa loi interne. La table multiplicative
nonadique occupe le corps principal ; chacune de ses lignes est engendrée par
accumulation additive. L'ombrage violet indique les lignes ou colonnes non
unitaires ; la ligne bleue identifie le sous-ensemble déterminé par le
représentant \(9\) de la classe nulle ; le cadre doré souligne le bloc
central \(2\times2\).}
\label{fig:app-ortograma}
\end{figure}

Chaque point de \(X_9\) reçoit simultanément les valeurs du produit et de la somme,
mais n'appartient pas pour autant à deux matrices APP distinctes. Les chemins de
\(\operatorname{Path}(\GraphAPP)\) parcourent un unique support ; \(\Pi\)
publie sa table principale et \(\Sigma\) conserve l'opération interne qui
engendre les accumulations et permet de mesurer leur incompatibilité.

\begin{proposition}[Structure des lignes et des colonnes]
\label{prop:app-latina-unidades}
La matrice de \(\Sigma\) est un carré latin d'ordre neuf. La matrice de
\(\Pi\) ne l'est pas ; cependant, pour chaque
\(u\in(\ZZ/9\ZZ)^\times=\{1,2,4,5,7,8\}\), la ligne et la colonne
correspondant à \(u\) sont des permutations des neuf classes résiduelles.
\end{proposition}

\begin{proof}
Une fois \(i\) fixé, l'application \(j\mapsto i+j\) est une translation de
\(\ZZ/9\ZZ\) ; il en va de même en fixant \(j\). Ainsi, chaque représentant
apparaît exactement une fois dans chaque ligne et chaque colonne de \(\Sigma\). Pour
\(\Pi\), la multiplication par \(u\) permute les classes si et seulement si \(u\)
est une unité. Les lignes correspondant à \(3,6,9\) présentent des répétitions,
de sorte que la propriété latine ne s'étend pas à toute la matrice
multiplicative.
\end{proof}

\section{Structure multiplicative de l'anneau local}

Chaque ligne de la table multiplicative est une addition répétée. La ligne \(i\)
accumule \(i\) le long d'une direction ; la colonne \(j\) accumule \(j\)
le long de la direction perpendiculaire. La cellule \((i,j)\) est l'intersection
des deux accumulations et renvoie \(ij=ji\). La table représente ainsi le produit
et la commutativité au moyen de deux accumulations transversales sur la même
grille.

La deuxième ligne est
\[
2,4,6,8,1,3,5,7,9.
\]
Avant le retour modulaire apparaissent les nombres pairs, et ensuite les impairs. La troisième
ligne est
\[
3,6,9,3,6,9,3,6,9,
\]
et rend le radical visible. Soit
\[
 R:=\ZZ/9\ZZ,
 \qquad
 \mathfrak m:=(3)=\{\overline0,\overline3,\overline6\}.
\]
La multiplication par trois définit
\[
 M_3:R\longrightarrow R,
 \qquad x\longmapsto3x,
 \qquad
 \ker M_3=\operatorname{im}M_3=\mathfrak m.
\]
Par le premier théorème d'isomorphisme, elle induit une identification linéaire
\begin{equation}
 \overline M_3:R/\mathfrak m\xrightarrow{\ \sim\ }\mathfrak m,
 \qquad \overline x\longmapsto3x,
 \label{eq:multiplicador-tres-radical-rev3}
\end{equation}
d'espaces vectoriels sur \(\mathbb F_3\). Ce n'est pas un isomorphisme
d'anneaux : \(\mathfrak m^2=0\), tandis que le quotient est un corps.

En représentants positifs, les trois phases de
\eqref{eq:multiplicador-tres-radical-rev3} sont
\begin{equation}
 369\longmapsto693\longmapsto936\longmapsto369,
 \label{eq:orbita-radical-369-rev3}
\end{equation}
produites par le déplacement de phase \(j\mapsto j+1\). La ligne
correspondant à \(6\equiv-3\pmod9\) inverse l'orientation du même
cycle. Cette orbite radicale ne doit pas être confondue avec l'entier
\(369\,693\,100\), qui apparaîtra ensuite comme nombre de chiffres de
\(9^{9^9}\).

En réduction directe modulo trois,
\[
3,6,9\longmapsto0,0,0.
\]
Si nous extrayons d'abord le facteur trois et conservons sa coordonnée interne,
\[
\eta_{\mathfrak m}(3a)=a\bmod3,
\qquad
3,6,9\longmapsto1,2,0.
\]
La seconde application conserve la coordonnée interne de l'idéal engendré par
trois. Elle ne doit pas être confondue avec la projection ternaire orientée qui sera définie
dans la section consacrée au TRIT\@.

\section{Sous-ensemble de coordonnées de la classe nulle}

Dans le graphe de l'observable multiplicative,
\[
\Pi(9,j)=\Pi(i,9)=9.
\]
La dernière ligne et la dernière colonne forment une union de dix-sept positions,
toutes de valeur neuf. Elles ne constituent pas une frontière intrinsèque du tore
arithmétique : c'est le sous-ensemble de la carte dans lequel au moins une coordonnée
représente la classe nulle par \(9\). Il existe en outre quatre représentants nuls intérieurs en
\[
(3,3),(3,6),(6,3),(6,6).
\]
La figure distingue ainsi ce sous-ensemble et les zéros internes de l'idéal
non unitaire. Le symbole supplémentaire employé plus loin dans la complétion
cubique appartient à un espace étendu et ne s'identifie à aucun
d'entre eux.

\section{Extension résidu--quotient}

La réduction modulo neuf conserve le résidu et élimine le quotient de la
division euclidienne. Pour étudier les deux composantes, nous définissons, en chaque point,
\[
i+j=9q^+_{ij}+R^+_{ij},
\qquad
ij=9q^\times_{ij}+R^\times_{ij},
\]
où \(R^+=\Sigma\), \(R^\times=\Pi\) et les \(q\) conservent le quotient.
L'extension résidu--quotient d'APP est notée
\[
\AAPPlift=(R^+,q^+;R^\times,q^\times).
\]

Soit \(s_9:\ZZ/9\ZZ\to\{1,\ldots,9\}\) la section de représentants fixée
par \(\dr\). La retenue additive associée à cette section est
\[
 \kappa_9(a,b)
 =\frac{s_9(a)+s_9(b)-s_9(a+b)}9\in\ZZ.
\]

\begin{lemma}[Identité cocyclique de la retenue]
\label{lem:app-cociclo-acarreo}
Pour tous \(a,b,c\in\ZZ/9\ZZ\),
\[
 \kappa_9(a,b)+\kappa_9(a+b,c)
 =\kappa_9(b,c)+\kappa_9(a,b+c).
\]
\end{lemma}

\begin{proof}
Les deux membres sont égaux à
\[
 \frac{s_9(a)+s_9(b)+s_9(c)-s_9(a+b+c)}9.\qedhere
\]
\end{proof}

Ce cocycle enregistre l'écart entre additionner des représentants entiers et
additionner d'abord dans le quotient. C'est la forme compositionnelle de la coordonnée de
quotient qui réapparaîtra dans la concaténation en base \(1000\).

Les totaux bruts sont
\[
\sum_{i,j}(i+j)
=9\sum_i i+9\sum_j j
=2\cdot9\cdot45
=810,
\]
et
\[
\sum_{i,j}ij
=\left(\sum_i i\right)\left(\sum_j j\right)
=45^2
=2025.
\]
Les sommes visibles s'obtiennent à partir des tables :
\[
\sum R^+=405,
\qquad
\sum R^\times=459.
\]
Par conséquent, les sommes totales des coordonnées de quotient sont
\[
\sum q^+=\frac{810-405}{9}=45,
\qquad
\sum q^\times=\frac{2025-459}{9}=174.
\]

\begin{exactbox}[title={Décomposition exacte de la différence intégrée}]
\[
\boxed{
2025-810
=1215
=(459-405)+9(174-45)
=54+9\cdot129.}
\]
La différence entre les sommes de résidus est \(54\), tandis que la
différence entre les sommes de quotients est \(129\). L'égalité recompose la
différence totale entre les observables entières et démontre que la composante
de quotient contient une information arithmétique non déterminée par la classe
résiduelle.
\end{exactbox}

\fi

\ifdefined\HMTAPPDevelopments
\chapter{Opérateurs arithmétiques et géométrie locale de l'espace des chemins}
\label{chap:app-operadores}

\section{Localisation de l'excès dans le radical nilpotent}

Le groupe des unités de \(\ZZ/9\ZZ\) est
\[
(\ZZ/9\ZZ)^\times=\{1,2,4,5,7,8\},
\]
et l'idéal maximal est
\[
\mathfrak m=(3)=\{0,3,6\}=\{9,3,6\}_{\rm visual},
\qquad\mathfrak m^2=0.
\]
Chaque ligne unitaire de \(\Pi\) est une permutation de \(1,\ldots,9\) et a pour somme
\(45\). Les lignes \(3\) et \(6\) ont pour somme \(54\) ; la ligne \(9\), \(81\).
Donc
\[
459=6\cdot45+2\cdot54+81
\]
et
\[
459-405=2(54-45)+(81-45)=54.
\]
Tout l'excès se localise dans le secteur non unitaire.

L'union des lignes et des colonnes \(3,6,9\) contient \(45\) cellules et a pour somme
\(297\). Son carré central \(3\times3\) a pour somme \(81=3\cdot27\) ; les bras
extérieurs, \(216=3\cdot72\). Ainsi,
\[
297=216+81=3(72+27)=3\cdot99.
\]
L'interprétation de \(72\) et de \(27\) s'obtient plus loin à partir du
bloc central, sélectionné de manière unique.

\section{Espérance conditionnelle par rapport à la projection \(\mathbb Z/9\mathbb Z\to\mathbb Z/3\mathbb Z\)}

La réduction canonique
\(\ZZ/9\ZZ\to\ZZ/3\ZZ\) détermine la partition
\[
 C_1=\{1,4,7\},\qquad
 C_2=\{2,5,8\},\qquad
 C_0=\{3,6,9\}.
\]
L'espérance conditionnelle uniforme moyenne chaque observable sur les neuf points de chaque bloc
\(C_a\times C_b\). Les matrices agrégées sont
\[
 M_{\Pi}=
 \begin{pmatrix}4&5&6\\5&4&6\\6&6&9\end{pmatrix},
 \qquad
 M_{\Sigma}=
 \begin{pmatrix}5&6&4\\6&4&5\\4&5&6\end{pmatrix}.
\]
Puisque chaque bloc contient neuf positions,
\[
 9\sum_{a,b}(M_\Pi)_{ab}=459,\qquad
 9\sum_{a,b}(M_\Sigma)_{ab}=405.
\]
Ainsi, le moyennage conserve les totaux des deux observables et localise de
nouveau leur différence :
\[
 \sum M_\Pi-\sum M_\Sigma=51-45=6,\qquad 9\cdot6=54.
\]

\begin{theorem}[Spectre de l'opérateur multiplicatif agrégé]
\label{thm:app-espectro-agregado}
L'opérateur matriciel \(M_\Pi\) satisfait
\[
 \det M_\Pi=-9
\]
et son spectre réel est
\[
 \operatorname{Spec}(M_\Pi)
 =\{-1,\,9-6\sqrt2,\,9+6\sqrt2\}.
\]
En particulier, la forme quadratique associée possède la signature \((2,1)\).
\end{theorem}

\begin{proof}
Le développement du déterminant donne \(-9\). Le polynôme caractéristique
se factorise sous la forme
\[
 (\lambda+1)\bigl((\lambda-9)^2-72\bigr).
\]
Ses racines sont celles indiquées et \(9-6\sqrt2>0\) ; il existe donc deux
directions positives et une négative.
\end{proof}

La décomposition spectrale admet une forme entière sur un sous-réseau
invariant. Pour
\[
 v_-=(1,-1,0),\qquad v_+=(1,1,0),\qquad v_0=(0,0,1),
\]
on a \(M_\Pi v_-=-v_-\), tandis que la restriction de \(M_\Pi\) au
sous-module engendré par \(v_+,v_0\) est représentée par
\[
 \begin{pmatrix}9&6\\12&9\end{pmatrix}
 =3H,\qquad
 H=\begin{pmatrix}3&2\\4&3\end{pmatrix}\in SL(2,\ZZ).
\]
Les valeurs propres de \(H\) sont
\[
 3\pm2\sqrt2=(1\pm\sqrt2)^2.
\]
Les trois vecteurs \(v_-,v_+,v_0\) engendrent un sous-réseau d'indice deux dans
\(\ZZ^3\), car le déterminant de la matrice qui les a pour colonnes vaut
\(-2\). On n'affirme donc pas ici une somme directe entière de tout
\(\ZZ^3\) : sur ce sous-réseau invariant, l'opérateur agrégé sépare une
direction de valeur propre \(-1\) et un sous-module hyperbolique régi par
l'unité quadratique \(3+2\sqrt2\) ; sur \(\RR^3\), la décomposition spectrale
est complète.

\section{Commutateur des projecteurs arithmétiques et séparation de leurs images}

La différence entre \(\Sigma\) et \(\Pi\) ne se réduit pas à leurs sommes globales.
Pour \(d\ge2\), soit \(\mathscr A_d=\{1,\ldots,9\}^d\) et définissons
\[
 \Sigma_d(x)=\dr(x_1+\cdots+x_d),\qquad
 \Pi_d(x)=\dr(x_1\cdots x_d).
\]
Dans l'espace de Hilbert \(\ell^2(\mathscr A_d)\), muni de la mesure
uniforme, notons \(P_{\Sigma,d}\) et \(P_{\Pi,d}\) les espérances
conditionnelles orthogonales sur les fonctions mesurables par rapport à
\(\Sigma_d\) et à \(\Pi_d\), respectivement. La table conjointe normalisée est
\[
 C^{(d)}_{ab}
 =\frac{N^{(d)}_{ab}}{\sqrt{n_a m_b}},
\quad
 N^{(d)}_{ab}
 =\#\{x:\Sigma_d(x)=a,\ \Pi_d(x)=b\},
\]
où \(n_a=\sum_bN^{(d)}_{ab}\) et \(m_b=\sum_aN^{(d)}_{ab}\).

Sur le support bidimensionnel original, la corrélation conjointe exacte est
\[
N^{(2)}=
\begin{pmatrix}
0&0&2&0&0&2&3&0&2\\
3&0&2&0&0&2&0&0&2\\
0&2&0&0&2&0&0&2&3\\
0&0&2&3&0&2&0&0&2\\
0&0&2&3&0&2&0&0&2\\
0&2&0&0&2&0&0&2&3\\
3&0&2&0&0&2&0&0&2\\
0&0&2&0&0&2&3&0&2\\
0&2&0&0&2&0&0&2&3
\end{pmatrix}.
\]
Les lignes, ordonnées selon la valeur de \(\Sigma\), ont pour somme neuf ; les colonnes,
ordonnées selon la valeur de \(\Pi\), ont pour somme
\((6,6,12,6,6,12,6,6,21)\). La corrélation conjointe des deux observables est ainsi déterminée par énumération complète
avant l'introduction de toute
application d'évaluation ultérieure.

\begin{proposition}[Commutateurs en dimensions deux et trois]
\label{prop:app-no-conmutatividad}
Les deux paires de projecteurs satisfont
\[
 \|[P_{\Sigma,2},P_{\Pi,2}]\|=\frac{\sqrt2}{3},
 \qquad
 \|[P_{\Sigma,3},P_{\Pi,3}]\|=\frac{\sqrt3}{4}.
\]
En particulier, les observables additive et multiplicative n'induisent pas de
décompositions simultanément compatibles.
\end{proposition}

\begin{proof}
Les valeurs singulières \(s\) de \(C^{(d)}\) sont les cosinus des angles
principaux entre les deux sous-espaces. Dans le bloc bidimensionnel
correspondant, la norme du commutateur est \(s\sqrt{1-s^2}\). L'énumération
exacte des \(9^d\) mots produit, pour \(d=2\), les valeurs propres non
triviales
\[
 s^2\in\left\{\frac57,\frac13,\frac13\right\}.
\]
Le maximum de \(s^2(1-s^2)\) est \(2/9\), d'où résulte
\(\sqrt2/3\). Pour \(d=3\), la même construction donne le maximum
\(s^2(1-s^2)=3/16\), atteint en \(s^2=1/4\), et la norme est donc
\(\sqrt3/4\). Dans les deux cas, le commutateur est non nul.
\end{proof}

\begin{proposition}[Intersection des sous-espaces observables]
\label{prop:app-interseccion-proyectores}
L'énumération exacte des mots de \(\mathscr A_d\), pour
\(2\leq d\leq15\), vérifie
\[
 \operatorname{Ran}P_{\Sigma,d}\cap
 \operatorname{Ran}P_{\Pi,d}
 =\CC\mathbf 1.
\]
Dans cet intervalle de dimensions, les fonctions constantes constituent
l'unique sous-espace commun aux deux décompositions arithmétiques.
\end{proposition}

\begin{proof}
Pour deux projections orthogonales, la dimension de l'intersection de leurs
images coïncide avec la multiplicité de la valeur singulière \(1\) dans la matrice
de recouvrement \(C^{(d)}\). Le calcul entier des tables
\(N^{(d)}\), suivi du calcul exact de cette multiplicité, donne la valeur un
pour chaque \(d=2,\ldots,15\). Le vecteur constant appartient aux deux images ;
il engendre donc toute l'intersection.
\end{proof}

La suite des normes ne s'extrapole pas à partir des deux premiers cas.
En particulier, pour \(d=4\), le même calcul donne
\[
 \|[P_{\Sigma,4},P_{\Pi,4}]\|
 =\sqrt{\frac{2618939}{22240656}}
 \simeq0.3431538652,
\]
ce qui ne coïncide pas avec la prolongation élémentaire \(2/(d+1)\). Ce cas
constitue un contre-exemple à toute formule inférée uniquement de
\(d=2,3\).

\section{Bloc central et décomposition spectrale}

La caractérisation du bloc central s'obtient au moyen d'une condition de
congruence. Pour deux représentants consécutifs \(k,k+1\), considérons
\[
B_k=
\begin{pmatrix}
k^2&k(k+1)\\
k(k+1)&(k+1)^2
\end{pmatrix}\pmod9.
\]

\begin{theorem}[Unicité du bloc diagonal adjacent]
\label{thm:app-bloque-central}
Il existe un unique bloc \(B_k\) dont les entrées diagonales coïncident. Il est
déterminé par \(k=4\) et vaut
\[
B_c=
\begin{pmatrix}7&2\\2&7\end{pmatrix}.
\]
\end{theorem}

\begin{proof}
L'égalité des diagonales équivaut à
\[
k^2\equiv(k+1)^2\pmod9,
\]
c'est-à-dire \(2k+1\equiv0\pmod9\). Comme \(2\) est une unité de
\(\ZZ/9\ZZ\), la congruence a pour unique solution \(k\equiv4\pmod9\).
\end{proof}

Par conséquent, les positions \(4,5\) de l'observable multiplicative forment
\[
B_c=
\begin{pmatrix}7&2\\2&7\end{pmatrix}
=7I+2R,
\qquad
R=\begin{pmatrix}0&1\\1&0\end{pmatrix},
\qquad R^2=I.
\]
Avec \(P_\pm=(I\pm R)/2\),
\[
B_c=9P_++5P_-.
\]
\begin{theorem}[Hiérarchie typée de la séparation somme--produit]
\label{thm:app-jerarquia-cuatro-dos-seis-cincuentaycuatro}
Le bloc central a un écart spectral local (9-5=4). Son coefficient
hors diagonale est (2). La compression par les trois classes modulo trois
produit la différence agrégée (51-45=6), et le déploiement sur les neuf
positions produit le défaut global
\[
 9\cdot6=459-405=54.
\]
Les entiers (4), (2), (6) et (54) appartiennent, respectivement, à
l'écart spectral, au couplage local, à la compression modulaire et à
l'intégration bidimensionnelle.
\end{theorem}

\begin{proof}
La décomposition \(B_c=7I+2R=9P_++5P_-\) donne l'écart et le
coefficient local. Les matrices agrégées \(M_\Pi,M_\Sigma\) construites
ci-dessus satisfont
\(\sum M_\Pi=51\) et \(\sum M_\Sigma=45\). Chaque entrée agrégée représente
neuf positions, de sorte que le déploiement multiplie la différence par
neuf et récupère les totaux (459) et (405).
\end{proof}
Leurs concaténations ordonnées produisent les blocs \(72\) et \(27\) :
\[
72+27=99,
\qquad
72-27=45=\det B_c.
\]
La concaténation des entrées diagonales et antidiagonales produit,
respectivement, \(77\) et \(22\). Par conséquent,
\[
 72+27=77+22=99,\qquad
 77-22=55=54+1.
\]
Ces égalités expriment deux partitions exactes du même bloc ; elles
n'identifient pas à elles seules une constante externe.

\begin{proposition}[Relation locale--globale du secteur radical]
\label{prop:app-local-global-radical}
La partition par lignes du bloc central reproduit, avec un facteur d'échelle
trois, la partition du secteur radical :
\[
 3\cdot72=216,\qquad
 3\cdot27=81,\qquad
 3\cdot99=297.
\]
\end{proposition}

\begin{proof}
Le bloc \(C_0\times C_0\) contient neuf représentants nuls et a pour somme
\(81\). Les quatre bras restants des lignes et des colonnes de \(C_0\)
ont pour somme \(216\). Les identités découlent de \(72+27=99\).
\end{proof}

De plus,
\[
M_c:=\frac{B_c}{\sqrt{45}}
=\exp\left(\frac12\log\frac95\,R\right)
\in SO^+(1,1).
\]
Si \(\eta=\operatorname{diag}(1,-1)\), la matrice normalisée \(M_c\) satisfait
\(M_c^{\mathsf T}\eta M_c=\eta\), \(\det M_c=1\) et appartient à
\(SO^+(1,1)\). Cette réalisation finie se limite au bloc \(2\times2\)
construit ; son interprétation spatio-temporelle requiert le pont dimensionnel
développé dans le traité de mécanique dimensionnelle.

\begin{proposition}[Famille spectrale APP--TRIT et lecture
Markov--Lorentz locale]
\label{prop:app-trit-espectral-markov-lorentz-rev8}
Pour \(\tau\in\{-1,0,+1\}\), soit
\[
 B_\tau=(4+3\tau)I+(5-3\tau)R.
\]
Alors
\[
 \operatorname{Spec}(B_\tau)=\{9,\,6\tau-1\},
 \qquad
 B_\tau=9P_++(6\tau-1)P_-,
\]
et le mode différentiel récupère l'orientation au moyen de
\[
 \tau=\frac{\lambda_-+1}{6}.
\]
En particulier, \(B_{+1}=B_c\). Sa normalisation stochastique
\[
 P_B:=\frac{1}{9}B_c=P_++\frac59P_-
\]
satisfait, pour tout \(n\geq0\),
\[
 P_B^n=P_++\left(\frac59\right)^n P_-.
\]
Par conséquent, la contraction du mode antisymétrique est
\(\kappa_B=5/9\) et l'écart spectral de cette chaîne locale est \(4/9\).
Si
\[
 \eta_c:=\frac12\log\frac95,
\]
la normalisation lorentzienne du même bloc satisfait
\[
 \frac{B_c}{\sqrt{45}}=\exp(\eta_c R),
 \qquad
 \boxed{\frac59=e^{-2\eta_c}}.
\]
\end{proposition}

\begin{proof}
Comme \(R\) agit avec les valeurs propres \(+1\) et \(-1\) sur
\(\operatorname{im}P_+\) et \(\operatorname{im}P_-\), respectivement, les
deux valeurs propres de \(B_\tau\) sont
\[
 (4+3\tau)+(5-3\tau)=9,\qquad
 (4+3\tau)-(5-3\tau)=6\tau-1.
\]
La formule des puissances découle de
\(P_\pm^2=P_\pm\) et de \(P_+P_-=0\). Enfin,
\(e^{-2\eta_c}=e^{-\log(9/5)}=5/9\).
\end{proof}

Cette coïncidence identifie deux lectures exactes du bloc local : la
mémoire différentielle décroît avec le même quotient que celui qui encode la rapidité
dans la réalisation déployée. Elle ne transforme pas \(B_\tau\) en le noyau de
Markov global du TPK, ni \(5/9\) en une constante physique. Elle n'étend pas non plus
à elle seule la réalisation \(SO^+(1,1)\) à un espace de Minkowski
\(3+1\) : ces étapes requièrent des applications dimensionnelles supplémentaires.

\section{Courbure discrète induite par l'interaction des feuillets somme et produit et son défaut de commutation}

Sur le tore, nous écrivons, maintenant en résidus,
\[
S(x,y)=x+y,
\qquad
P(x,y)=xy\pmod9.
\]
Pour toute fonction \(T:X_9\to\ZZ/9\ZZ\), nous fixons les différences progressives
\[
 \Delta_xT(x,y)=T(x+1,y)-T(x,y),\qquad
 \Delta_yT(x,y)=T(x,y+1)-T(x,y),
\]
où les coordonnées sont prises modulo neuf. Cette convention fixe également
l'orientation positive de chaque plaquette : on parcourt d'abord la direction
\(x\), puis la direction \(y\), et l'on revient par les arêtes opposées.
Il est important de typer correctement ces deux objets. \(S\) et \(P\) sont
des potentiels de sommet ; les variables de lien sont leurs différences
discrètes
\[
A_x^T=\Delta_xT,
\qquad
A_y^T=\Delta_yT.
\]
Pour un mélange ordonné de potentiels \((T_x,T_y)\), nous définissons la
courbure de plaquette par
\[
F(T_x,T_y)
=\Delta_x A_y^{T_y}-\Delta_y A_x^{T_x}
=\Delta_x\Delta_yT_y-\Delta_y\Delta_xT_x.
\]
Comme
\[
\Delta_xS=\Delta_yS=1,
\qquad
\Delta_xP=y,
\qquad
\Delta_yP=x,
\]
on obtient sur les \(81\) plaquettes :
\[
F(S,S)=F(P,P)=0,
\qquad
F(S,P)=+1,
\qquad
F(P,S)=-1\pmod9.
\]
Les connexions associées à une unique observable sont plates ; la composition
mixte des deux observables produit une courbure de signe opposé selon
l'ordre. Pour un rectangle de coordonnées \(R_{a,b}\) contenu
dans la carte, nous définissons le Wilson additif de bord comme la somme orientée des
variables de lien. L'annulation des arêtes intérieures donne
\[
 W_{\partial R_{a,b}}(S,P)=ab,\qquad
 W_{\partial R_{a,b}}(P,S)=-ab\pmod9,
\]
qui est l'identité de Stokes finie pour cette convention.

La distinction entre potentiel et lien est algébriquement nécessaire. Si l'on remplaçait
directement \(A_x,A_y\) par \(S,P\), des termes dépendant de
\(x,y\) apparaîtraient et les valeurs constantes précédentes ne seraient pas obtenues. La formulation
typée est celle qui produit le cocycle mixte utilisé plus loin dans la fermeture
de Heisenberg finie.

\begin{remark}
La différence intégrée \(54\) et la courbure mixte \(\pm1\) sont des invariants
distincts. La première compare des sommes globales ; la seconde dépend de l'ordre des
deux potentiels et se localise sur chaque plaquette. La formulation au moyen
d'opérateurs autoadjoints et de commutateurs sera développée après la construction de
l'espace d'états correspondant.
\end{remark}

\section{Développement de l'opérateur de transition par chemins}

La définition d'APP comprend l'espace
\(\operatorname{Path}(\GraphAPP)\). Une fois fixé un opérateur de
transition \(U\) sur les sommets de \(\GraphAPP\), nous supposons que
\(U_{yx}=0\) dès que \(x\to y\) n'est pas une arête du graphe. Sous cette
hypothèse de support, ses puissances admettent le développement
\begin{proposition}[Développement fini par chemins]
\label{prop:app-suma-caminos}
Pour chaque \(R\ge1\) et chaque paire de sommets \(i,f\),
\[
(U^R)_{fi}
=\sum_{\substack{\gamma:i\to f\\|\gamma|=R}}
\prod_{k=0}^{R-1}U_{x_{k+1},x_k},
\qquad
\gamma=(x_0=i,x_1,\ldots,x_R=f).
\]
\end{proposition}

\begin{proof}
Pour \(R=1\), l'identité est la définition des éléments de matrice de
\(U\). Si elle est vraie pour \(R\), alors
\[
 (U^{R+1})_{fi}
 =\sum_j U_{fj}(U^R)_{ji}.
\]
Substituer l'hypothèse de récurrence concatène à chaque chemin de longueur \(R\)
de \(i\) à \(j\) la dernière arête \(j\to f\) ; ainsi, chaque chemin de
longueur~\(R+1\) apparaît exactement une fois.\qedhere
\end{proof}

Par conséquent, APP réunit en une seule structure le tore arithmétique fini,
les observables de somme et de produit, la décomposition résidu--quotient, le
radical nilpotent et une courbure mixte orientée. Le représentant \(9\) de
la classe nulle fixe une carte visible du quotient ; le remplacer par \(0\) sans
transporter simultanément la carte modifie les registres utilisés par la
dynamique ultérieure.
\fi
